\documentclass[a4paper,12pt,french,twoside,footnotereset=true,versetitle=true]{droit-fr}
\usepackage[T1]{fontenc}
\usepackage{microtype}
\usepackage{hyperref}
% The CTAN droit-fr bibliography style is from 2016 and is incompatible with
% the pinned BibLaTeX version. Keep the legal layout and use maintained style.
\usepackage[style=numeric,backend=biber,indexing=cite]{biblatex}
\addbibresource{references.bib}
\hypersetup{pdftitle={Une étude sur la preuve numérique},pdfauthor={Claire Martin}}

\author{Claire}{Martin}
\title{La preuve numérique et la procédure contradictoire}
\university{de la République}
\school{Droit privé et sciences du numérique}
\speciality{Droit du numérique}
\approvaldate{29 août 2026}
\director{Madame}{Nora}{Leblanc}{Professeure à l'Université de la République}
\reportera{Monsieur}{Émile}{Durand}{Professeur à l'Université du Littoral}
\reporterb{Madame}{Inès}{Moreau}{Maîtresse de conférences à l'Université de la République}
\membera{Monsieur}{Paul}{Roux}{Avocat au barreau de Paris}
\memberb{Madame}{Sara}{Petit}{Magistrate honoraire}

\begin{document}
\maketitlepage

\frontmatter
\chapter*{Résumé}
Cette étude présente une grille de lecture pour les éléments de preuve
numérique lorsqu'une partie doit pouvoir discuter leur origine, leur intégrité
et leur interprétation. L'exemple est fictif et sert uniquement de point de
départ pour un mémoire en droit français.
\tableofcontents

\mainmatter
\partie{Cadre de l'étude}
\verset{}
La preuve numérique circule entre plusieurs acteurs et plusieurs systèmes.
L'analyse juridique doit donc conserver la chaîne de conservation, les
conditions de collecte et la possibilité d'une discussion contradictoire.

\titre{La chaîne de conservation}
\verset{}
Une copie horodatée n'est pas une explication complète de la preuve. Il faut
également indiquer qui a effectué la copie, avec quel outil, sur quel support,
et quelles opérations ont suivi.

\chapitre{Méthode de qualification}
\verset{}
Nous distinguons trois questions : l'identification de la source, l'intégrité
du fichier et la pertinence de l'élément pour le litige. Cette décomposition
permet d'éviter de confondre authenticité technique et force probante.

\section{Source et intégrité}
\verset{}
Le procès-verbal devrait préciser le contexte d'obtention et permettre à la
partie adverse de demander une vérification. L'horodatage, le calcul d'une
empreinte et la conservation du support original sont des éléments utiles,
mais chacun doit être relié à une personne et à une procédure identifiables.
\footnote{Dans une étude réelle, remplacer cet exemple par l'autorité et la
référence précises qui s'appliquent au dossier.}

\paragraphe{Une méthode reproductible}
\verset{}
Un protocole court peut lister les commandes utilisées, conserver leurs
sorties et joindre l'empreinte du fichier traité. Cette discipline technique
facilite ensuite le débat juridique sans prétendre résoudre seule la question
de la valeur probante.

\partie{Discussion et conclusion}
\verset{}
La preuve numérique ne devrait pas être présentée comme un objet autosuffisant.
La partie qui l'invoque doit exposer son parcours et laisser une possibilité
réelle de contrôle. Cette exigence protège à la fois l'efficacité de la preuve
et le principe du contradictoire.

\printbibliography
\end{document}
